\subsection{原生适配与独立接收的验证层次}
\label{sec:native-adapter-validation}
任务执行、证据归档和统计解释分别接受核验。独立冻结的原生适配验收在G1、G2和水井W1各使用一份实际输入，运行@@MAIN@@项四链主任务及@@CACHE@@项缓存探测；G1/W1每链保留64步，G2保留16384步。接受原生进程、故障隔离、恢复和身份拒绝检查后，主任务全部返回数值合格结果，缓存@@CALLS@@次调用仍无普通后验样本资格。前期句柄及测试环境失败、失败测试尝试与修复记录一并归档，没有将测试重试计为新的科学拟合。

接收端在Mac上按完整任务表重读原始证据，对@@MH@@项MH主工作流及全部缓存调用独立重放NumPy递推，接受事件失配均为@@MISMATCH@@；三项CPU NUTS按数组、随机状态、适应记录和子进程证据核对，不作MH同路径声明。随后从原始函数二进制重建全部@@DIAG@@条R诊断，按既定跨平台输出条件通过比较；这与早期要求原尺度输出逐位相等的检查是不同契约，没有将此前60份严格变换失败改为通过。恢复分析复用@@REUSED@@项结果，未重新采样或重算诊断，原始文件不变。

数值合格仍不能解释为后验探索充分：本次@@DIAG@@条函数记录中，@@BAD@@条有限$\hat R$大于1.01，@@NULL@@条$\hat R$及@@TAIL@@条tail ESS不可判定。这些记录共享每个目标的一份输入，不能据工作流或缓存调用数建立正式独立重复区间。逐函数诊断、跨平台末位差、数值库警告及原生故障链见配套接收报告；当前证据支持该固定适配版本的有限运行与接收能力，九目标多预算的正式推断比较仍未冻结或执行。
